\section{A split sixfold and the coniveau question for a Mumford square}
\label{sec:mumford-coniveau}

This section gives a second sufficient geometric condition for the
exceptional classes of a Mumford square, using a polarized abelian
sixfold and the first cohomology of a resolved divisor.
It makes precise the distinction between an algebraic Weil class and
the support assertion in the generalized Hodge conjecture.

Let $F$ be a totally real cubic field, let $D/F$ be a quaternion algebra
split at exactly one real place, and assume
\begin{equation}
 E=F(\sqrt{-3})\text{ splits }D,
 \qquad \operatorname{Cor}_{F/\mathbf Q}(D)\text{ is split}.
 \label{eq:mumford-sixfold-arithmetic}
\end{equation}
These conditions hold when $D$ is unramified at every finite place:
the quadratic extension is complex at the two ramified infinite places,
and the two local corestriction invariants sum to zero in
$\mathbf Q/\mathbf Z$.  In particular they include the quaternion
data of the two curve families in \cite[Theorem~1.2]{BHPS}.
Write $K=\mathbf Q(\sqrt{-3})$.  Let $A$ be the Mumford fourfold
with $V=H^1(A,\mathbf Q)$ and $V_{\mathbf C}=(2,2,2)$.

For a smooth projective variety $X$, let $N^1H^j(X,\mathbf Q)$ be the
sum of the kernels of restriction to $X\setminus Z$, where $Z$ ranges
over divisors.  This is geometric coniveau one; the assertion that a
Hodge substructure has no $(j,0)$ or $(0,j)$ part does not by itself
put it in this subspace.

\begin{theorem}[A coniveau-one criterion]\label{thm:mumford-coniveau}
Under \eqref{eq:mumford-sixfold-arithmetic}, there is a polarized split
Weil sixfold $B$ for $K$, varying over a common finite cover of the
Mumford Shimura curve, and an algebraic cohomological projector with
image
\[
 P\subset H^3(B,\mathbf Q),\qquad \dim_{\mathbf Q}P=16,
 \qquad P(1)\simeq V^{\oplus2}.
\]
At a Hodge-generic point of the common $(A,B)$ family, the condition
\begin{equation}
 P\cap N^1H^3(B,\mathbf Q)\ne0
 \label{eq:mumford-coniveau-input}
\end{equation}
implies the rational Hodge conjecture for every power of $A$.
In particular, the whole exceptional plane in $H^4(A^2,\mathbf Q)$ is
algebraic.  Equivalently, it suffices to have one nonzero algebraic
correspondence $H^1(A)\longrightarrow P(1)$.
The theorem does not assert \eqref{eq:mumford-coniveau-input}.
\end{theorem}

The proof first constructs $B$ and its projector using endomorphisms.
A supporting divisor would then give an Albanese factor isogenous to
$A$.  Divisor invariants for the larger group of $E$-linear unitary
transformations on $H^1(B)$ transfer through this factor to all
Mumford invariants.  We also compute the resulting weighted Casimir
on the exceptional plane explicitly.
Bhattacharjee--Bhattacharya \cite[Theorems~37--38]{BBK3Sixfolds}
give conditional transfer criteria using an associated sixfold and
an adjoint K3 structure.  Here the sixfold has the balanced compact
signatures $(2,0)$ and $(0,2)$; the required input is precisely
geometric coniveau of its rank-sixteen degree-three summand.

\subsection{The Hermitian space and its determinant}

Since $E$ splits $D$, write $D=(-3,b)_F$, with $b\in F^\times$;
the quadratic splitting criterion is
\cite[Lemma~5.4.7, pp.~73--74]{VoightQuaternion}.
Label the real embeddings so that the signs of $b$ are $(+,-,-)$.
Choose $\delta\in F^\times$ of signs $(+,-,+)$ and set
\[
 h_E=\delta\begin{pmatrix}b&0\\0&-1\end{pmatrix},
 \qquad h_K=\operatorname{Tr}_{E/K}h_E.
\]
The three $E/F$ signatures of $h_E$ are $(1,1),(2,0),(0,2)$;
its special unitary group is $\operatorname{SL}_1(D)$.
The $K$-Hermitian space $(E^2,h_K)$ has dimension six and signature
$(3,3)$.

Let $d_F$ be the discriminant of a rational basis of $F$.  Taking the
determinants of the two trace-form blocks gives
\begin{equation}
 \det_K h_K
 =-\operatorname{Nm}_{F/\mathbf Q}(b)
       \bigl(\operatorname{Nm}_{F/\mathbf Q}(\delta)d_F\bigr)^2.
 \label{eq:mumford-sixfold-determinant}
\end{equation}
The projection formula for quaternion symbols
\cite[II, Proposition~1.39(d), p.~74]{MilneCFT} gives
\[
 \operatorname{Cor}_{F/\mathbf Q}(-3,b)_F
   =(-3,\operatorname{Nm}_{F/\mathbf Q}(b))_{\mathbf Q}.
\]
Thus the first norm in \eqref{eq:mumford-sixfold-determinant} is a
$K/\mathbf Q$ norm, by the split-corestriction hypothesis and
\cite[Main Theorem~5.4.4(vi), pp.~72--73]{VoightQuaternion}.  The
determinant class is consequently $-1$, the class of three hyperbolic
planes.  Landherr's classification by dimension, determinant and
signatures \cite[pp.~246--247]{Landherr} identifies $h_K$ with that
hyperbolic space.

On the noncompact real factor choose the varying Hodge line of the
Mumford disk.  On the two compact factors choose respectively all of
the indicated $K$ eigenspace and none of it as the $(1,0)$ part.
Tracing $h_E$ against a purely imaginary element of $K$, with the
corresponding sign, gives a positive alternating Riemann form.
An $E$-stable lattice and an integral multiple of this form give the
polarized abelian sixfold $B$.  Its $K$ action has signature $(3,3)$,
and the preceding determinant calculation proves that it is split.
This construction does not require a principal polarization.
Fixed lattices and sufficiently fine level realize the two rational
representations as abelian schemes over a common finite cover of the
Mumford curve; the compact central Hodge data are fixed throughout.

\subsection{An algebraic projector onto the rank-sixteen part}

Put $H=H^1(B,\mathbf Q)$.  For one embedding $\sigma:K\to\mathbf C$,
write
\[
 H_\sigma=M_1\oplus M_2\oplus M_3,\qquad \dim M_i=2.
\]
The three Hodge-circle weight lists are $(1,-1),(1,1),(-1,-1)$.
The multidegree $(1,1,1)$ part of $\Lambda_K^3H$, together with its
conjugate, is
\begin{equation}
 P_{\mathbf C}=(M_1\otimes M_2\otimes M_3)
       \oplus(\overline M_1\otimes\overline M_2\otimes\overline M_3).
 \label{eq:mumford-rank16}
\end{equation}
It is Galois stable.  Its Hodge types are $(2,1)$ and $(1,2)$, each
with dimension eight.  The connected central torus acts trivially on
it: the Galois conjugates of the central cocharacter $(0,1,-1)$ span
the trace-zero plane, whereas the tensor in each summand has the
product character.  The derived representation is $(2,2,2)$.
The split corestriction gives its rational eight-dimensional form
$V$, with endomorphism field $\mathbf Q$.  The Hodge circle just
computed therefore identifies $P(1)$ with two copies of $V$.
These are statements at a Hodge-generic point of the common family;
they do not use the enlarged endomorphism algebra of a CM fibre.

We now give a projector without assuming that a general Hodge
projector is algebraic.  Algebraic K\"unneth projectors on an abelian
variety are polynomials in a multiplication map.  On $H^3(B)$ first use
\[
 e_K=\frac{1+[\omega]^*+[\omega^2]^*}{3}.
\]
The eigenvalue of $[\omega]^*$ on a wedge with $r$ factors from
$H_\sigma$ is $\omega^{2r-3}$.  Consequently $e_K$ selects exactly
$\Lambda^3H_\sigma\oplus\Lambda^3H_{\bar\sigma}$.

Choose a primitive element $a\in F$ with conjugates $a_1,a_2,a_3$.
Let $D_a$ apply $a$ in one of the three exterior factors and the
identity in the others, summed over the three choices.  It is an
algebraic cohomological operator, since
\begin{equation}
 D_a=\frac{[a+2]^*-2[a+1]^*+[a]^*}{2}-3\operatorname{id}
       \quad\hbox{on }H^3(B).
 \label{eq:mumford-derivation}
\end{equation}
Indeed, the second finite difference of $\Lambda^3(a+t)$ is
$6\operatorname{id}+2D_a$.  A rational endomorphism is interpreted
using a rational graph correspondence, or after replacing the lattice
by an isogenous one on which it is integral.

On \eqref{eq:mumford-rank16} the eigenvalue is
$c=\operatorname{Tr}_{F/\mathbf Q}(a)$.
All the other multidegrees are $(2,1,0)$, with eigenvalues
$2a_i+a_j$, $i\ne j$.  Set
\[
 Q_a(T)=\prod_{i\ne j}(T-2a_i-a_j)\in\mathbf Q[T].
\]
Then
\[
 Q_a(c)=-\prod_{i<j}(a_i-a_j)^2\ne0,
 \qquad e_P=e_K\frac{Q_a(D_a)}{Q_a(c)}
\]
is the required algebraic projector.  Repeated unwanted eigenvalues
cause no problem.  This also constructs the projector relatively,
using fixed endomorphisms and rational coefficients.

\subsection{A divisor support gives the weight-one correspondence}

We explain precisely the support assertion used in
\eqref{eq:mumford-coniveau-input}.  The coniveau-one part of $H^3(B)$
is the sum of Gysin images
\begin{equation}
 H^1(\widetilde D,\mathbf Q)(-1)\longrightarrow H^3(B,\mathbf Q)
 \label{eq:mumford-divisor-gysin}
\end{equation}
from resolutions of divisors.  Equivalently it is the image of a
correspondence from the first cohomology of a smooth projective curve
\cite[Proposition~1.3, p.~8]{ACMV}.  For degree three this can also be seen
directly: remove the singular loci and intersections of the supporting
divisors.  They have codimension at least two and contribute no
cohomology with supports in degree three.  Purity for the remaining
smooth divisors gives their $H^1(-1)$.  Only its weight-one part can
map to the pure group $H^3(B)$, and this part comes from a smooth
projective resolution.  A finite sum suffices for any given class.

Apply $e_P$ to these maps.  Under
\eqref{eq:mumford-coniveau-input} at least one resulting map is
nonzero.  Polarizable rational Hodge structures are semisimple,
and $P(1)=V^{\oplus2}$ with $V$ simple.  Therefore the Albanese
variety of a product of the resolutions has an isogeny factor
isogenous to $A$.  Hodge morphisms between the first cohomology of
abelian varieties are induced by homomorphisms up to isogeny.
Composing such a map with \eqref{eq:mumford-divisor-gysin} and $e_P$
produces a nonzero algebraic correspondence
\begin{equation}
 \Gamma:V\longrightarrow P(1).
 \label{eq:mumford-coniveau-bridge}
\end{equation}

In the convention of a map from $A$ to $B$, this correspondence has
codimension five in $A\times B$: its degree shift is $3-1=2$.
Conversely, a codimension-five cycle on this ten-dimensional product
has dimension five.  Its image in the sixfold $B$ is contained in a
divisor, so its cohomological image is in $N^1H^3(B)$.
Degree-zero algebraic correspondences preserve coniveau, by moving
intersections properly and the support dimension inequality; thus
applying $e_P$ preserves this conclusion.

For comparison with a Hodge tensor, the relevant component of the
codimension-five cycle is in $H^7(A)\otimes H^3(B)$.
The inverse of $L_A^3:H^1(A)\to H^7(A)$ is algebraic.  In the notation
of Section~\ref{sec:mumford}, it is $\Lambda_A^3/36$, since
$\Lambda_A^3L_A^3=36$ on $H^1(A)$.
Applying it in the first factor gives a codimension-two cycle with
component in $H^1(A)\otimes H^3(B)$.  Multiplication by $\theta_A^3$
reverses this operation.  Neither description asserts that the bridge
is a divisor on $A\times B$.

\subsection{Divisor invariants transfer to every power}

Let $U$ be the rational algebraic group of $E$-linear automorphisms of
$H^1(B)$ preserving its alternating polarization.  It contains the
Hodge circle of $B$, and
\[
 U_{\mathbf C}\simeq\mathrm{GL}_2^3,
 \qquad H_{\mathbf C}=\bigoplus_{i=1}^3(M_i\oplus M_i^\vee).
\]
The indicated dualities include the polarization Tate twist.
The invariants of $U$ in the cohomology of powers of $B$ are generated
by divisors.  Indeed the first fundamental theorem for the general
linear group expresses every invariant tensor as a linear combination
of contractions between $M_i$ and $M_i^\vee$.  These contractions are
the polarization pairings with the $E$-action inserted.  They are
divisor classes on $B^2$; permutations, exterior products and diagonal
pullbacks give their images in all exterior powers.  The rational
descent and the divisor formulation are
\cite[Proposition~3.6(c) and Remark~3.7, pp.~17--18]{MilneLefschetz}.
This assertion concerns the larger group $U$.  It does not say that
all Hodge classes on $B$ are generated by divisors.

The space $P$ is primitive in $H^3(B)$: the lowering operator pairs
opposite $K$ eigenspaces, whereas each summand in
\eqref{eq:mumford-rank16} uses just one eigenspace.
Consequently the polarization adjoint of $\Gamma$ is algebraic.
More explicitly, on the relevant cohomological degrees it is
\[
 \beta=(L_A^3)^{-1}\,\Gamma^t\,L_B^3 e_P:P(1)\longrightarrow V,
\]
up to the fixed sign in the polarization convention.  Include the
algebraic degree-one projector on $A$.  The resulting correspondence
lies in $\mathrm{CH}^5(B\times A)_{\mathbf Q}$.
It is nonzero, since it is the adjoint of a nonzero map between
polarized Hodge structures.

Put $P_+=M_1\otimes M_2\otimes M_3$ and
$P_-=M_1^\vee\otimes M_2^\vee\otimes M_3^\vee$ over $\mathbf C$.
For $G_{\mathbf C}=\mathrm{SL}_2^3$, both are irreducible of the
same type as $V_{\mathbf C}$.  Each restriction $\beta|_{P_\pm}$
is an isomorphism of $G_{\mathbf C}$-modules: it is either zero or
an isomorphism, and rationality interchanges the two restrictions.
They cannot both be zero.

For an integer $m\geq1$ consider
\[
 I_m=(P^{\otimes 2m})^U.
\]
After complexification its summand with $m$ factors $P_+$ and $m$
factors $P_-$ contains all the $G_{\mathbf C}$-invariants of that
balanced tensor product.  To check this, regroup its factors as
\[
 \bigotimes_{i=1}^3
       \bigl(M_i^{\otimes m}\otimes(M_i^\vee)^{\otimes m}\bigr).
\]
The scalar centre of each $\mathrm{GL}_2$ acts trivially here.
Thus its $\mathrm{GL}_2$-invariants and
$\mathrm{SL}_2$-invariants coincide.
The isomorphisms $\beta|_{P_\pm}$ now show that
\begin{equation}
 \beta^{\otimes2m}:I_m\longrightarrow(V^{\otimes2m})^G
       \quad\hbox{is surjective}.
 \label{eq:mumford-balanced-invariants}
\end{equation}
Surjectivity over $\mathbf C$ implies surjectivity over $\mathbf Q$.

Every element of $I_m$ is a divisor-product class of codimension $3m$
on $B^{2m}$.  Applying $\beta^{\boxtimes2m}$ produces a cycle of
codimension $m$ on $A^{2m}$: the correspondence has codimension $10m$
and the source has dimension $12m$.  This proves algebraicity of
every tensor in the target of \eqref{eq:mumford-balanced-invariants}.
Finally, $H^*(A^r)=\Lambda^*(V^{\oplus r})$.  Exterior multiplication
is equivariant, and invariants are exact for the reductive group $G$.
Every invariant in degree $2m$ is therefore a sum of images of
invariant $2m$-fold tensors.  Assigning each tensor factor to one of
the $r$ coordinates is the pullback along a map
$A^r\to A^{2m}$ whose coordinates repeat the corresponding
projections.  Such pullbacks preserve algebraicity.
This proves the rational Hodge conjecture on every $A^r$, and hence
Theorem~\ref{thm:mumford-coniveau}.

\subsection{Transporting the Lie action and its field weights}

We give a direct description in degree four of the exceptional classes
obtained from the same bridge.

Inside $\operatorname{End}_{\mathbf Q}(H)$ consider
\[
 \mathfrak g=\{X\in\operatorname{End}_E(H):
                    X^*=-X,\ \operatorname{Tr}_E X=0\},
\]
where $*$ is the polarization adjoint.  It is the rank-nine
representation $\operatorname{Res}_{F/\mathbf Q}\mathfrak{sl}_1(D)$.
Its projector is algebraic for an explicit reason.  The separability
idempotent of $E/\mathbf Q$, expressed in trace-dual bases, projects
to $\operatorname{End}_E(H)$.  The operation $(1-*)/2$ takes its
skew-adjoint part.  Subtracting the scalar skew-adjoint $E$ part by
the rational trace pairing takes the trace-zero part.
All $E$ actions are endomorphisms of $B$; the adjoint, pairings and
their inverses use the polarization and the algebraic Lefschetz
operators.  Thus these operations are algebraic on powers of $B$.
Multiplication by any $a\in F$ on $\mathfrak g$ is algebraic too.

The derivation action on $\Lambda^3H$, followed by $e_P$, gives an
algebraic map $\mathfrak g\to\operatorname{End}(P(1))$.
At the common Hodge-generic point, \eqref{eq:mumford-coniveau-bridge}
intertwines this action with the Mumford action.  Its polarization
adjoint is algebraic and
\[
 \Gamma^\dagger\Gamma=r\operatorname{id}_V,
 \qquad 0<r\in\mathbf Q,
\]
because $\operatorname{End}_{\mathrm{HS}}(V)=\mathbf Q$.
Consequently
\[
 X\longmapsto r^{-1}\Gamma^\dagger\rho_P(X)\Gamma
\]
is the usual $\mathfrak g$ action on $V$, expressed by algebraic
correspondences.

Fix the trace form $B_{\mathfrak g}(X,Y)=\operatorname{Tr}_{\mathbf Q}(XY)$
in the $E^2$ representation.  It is nondegenerate and has the same
rational index on the three complex $\mathfrak{sl}_2$ factors.
Its inverse tensor is algebraic, using the preceding projector and
the ambient trace coevaluation.  After fixing this inverse, apply
multiplication by a nonrational $a\in F$ in one factor and transport
both factors to $\operatorname{End}(V)$.  The resulting tensor is,
up to a common nonzero rational constant, the weighted Casimir
$\sum_i a_i C_i$.  The separate multiplication by $a$ is essential:
it is not a change of metric followed by an inverse that cancels it.

Regroup its four $V$ factors into two exterior pairs.  The calculation
\eqref{eq:mumford-casimir-eigen} gives the three eigenvalues
\[
 \frac14\operatorname{Tr}_{F/\mathbf Q}(a)-a_j.
\]
For clarity, the cross Casimir is $1/4$ on the symmetric square of
the two-dimensional standard representation and $-3/4$ on its
exterior square.  In the $j$th primitive summand of $\Lambda^2V$,
the exterior choice occurs at $j$ and the symmetric choices at the
other two indices.  This gives exactly the displayed expression.
Algebraic primitive projection and trace subtraction leave
\[
 -a+\frac13\operatorname{Tr}_{F/\mathbf Q}(a)\in F^0\setminus\{0\}.
\]
The cubic-field composition argument of
Theorem~\ref{thm:mumford-cubic} now gives the entire exceptional plane.
This describes the exceptional plane explicitly.

\subsection{What the algebraic Weil classes do give}

The ordinary Weil classes on $B$ are algebraic by the split-sixfold
results, with the input status in Section~\ref{sec:effective-sixfolds}.
They provide the following intermediate conclusion, which does not
settle the support condition above.

\begin{proposition}\label{prop:weil-support-middle}
If a nonzero rational pure Weil class $w\in H^6(B,\mathbf Q)$ is
represented by a rational combination of integral threefolds $Y_i$,
then $P$ is an algebraic direct summand of
$\bigoplus_i H^3(\widetilde Y_i,\mathbf Q)$, where the tildes are
smooth projective resolutions.
\end{proposition}

\begin{proof}
For $i:\widetilde Y_i\to B$, the cup action of its cycle class is
$i_*i^*:H^3(B)\to H^9(B)$.
The inverse of $L_B^3:H^3(B)\to H^9(B)$ is algebraic, by the
Lefschetz lowering operator and its finite primitive decomposition.
Consider on $P$ the operator
\[
 T=e_P(L_B^3)^{-1}(w\smile -)e_P.
\]
It factors algebraically through the displayed sum of middle
cohomology groups.  We show it is invertible.

Over $\mathbf C$ write $H=H_\sigma\oplus H_{\bar\sigma}$.
The polarization pairs these spaces perfectly and respects their
three $F$ blocks.  The two components of $w$ are nonzero volume
forms on these six-dimensional spaces: rationality and $w\ne0$
make both coefficients nonzero.  For
$u\in M_1\otimes M_2\otimes M_3\subset\Lambda^3H_\sigma$,
only the conjugate volume contributes to $w\smile u$.
The restriction of $L_B^3$ from the conjugate triple wedge to this
bidegree is an isomorphism, given by complementary exterior-power
pairing.  In multidegree $(1,1,1)$ it uses one Lefschetz pair from
each $F$ block.  Thus $T$ exchanges the two summands of
\eqref{eq:mumford-rank16} by invertible maps.  Here the two volume
forms refer to $K$ eigenspaces; their Hodge type is $(3,3)$.

The inverse of $T$ on $P$ is a polynomial in $T$ with rational
coefficients, by its characteristic polynomial and nonzero
determinant.  Composing this inverse with the factorization gives
an algebraic splitting, as asserted.
\end{proof}

The proposition reaches $H^3$ of threefolds.  A coniveau-one class
on one such $\widetilde Y_i$ whose image under
$e_P(L_B^3)^{-1}i_*$ is nonzero would suffice: its Gysin
pushforward raises coniveau by three, and the inverse $L_B^3$
lowers it by at most three, leaving the additional step of
coniveau one on $B$.  Ordinary algebraicity of the degree-six
Weil class does not prove this generalized Hodge assertion.

There is also a direct limitation on translations.  At the generic
point, every $H^1$ of a power of $B$ has derived constituents
$(2,1,1),(1,2,1),(1,1,2)$; a CM parameter variety adds only
derived-trivial constituents.  Hence
\[
 \operatorname{Hom}_{\mathrm{HS}}
   (H^1(B^m\times C_{\mathrm{CM}}),P(1))=0.
\]
Any construction using these parameter first-cohomology groups,
addition, translations and Lefschetz correspondences therefore
has zero component in the required space.  This does not exclude
a singular divisor whose resolution has a new Albanese factor
isogenous to $A$; finding such a factor is the unresolved geometry.

Abdulali's type-IV theorem \cite[Theorem~13]{AbdulaliIV} requires
$|p_\alpha-q_\alpha|=1$ at every real place.  Here those differences
are $0,2,2$.  His other criterion requires every real simple factor
to be $\mathrm{SU}(p,1)$ with $p>1$
\cite[Theorems~6.5--6.6]{AbdulaliSurvey}; the compact factors here
do not satisfy it.  Appendix~B.5 of that survey adds the former
signature hypothesis to the earlier type-IV statements.
Its Section~8.2 records Schoen's generalized Hodge result for the
rank-two determinant structure of a signature-$(5,1)$ Eisenstein
sixfold.  That is a CM structure in degree six, not the rank-sixteen
part of degree three used here.  None of these cited statements
establishes \eqref{eq:mumford-coniveau-input}.

Finally, the family conclusion needs a family input.  A relative
nonzero correspondence \eqref{eq:mumford-coniveau-bridge} over a
nonempty open subset gives a uniform cycle bound after spreading
out, and proper bounded-cycle spaces extend it to the whole base.
The same follows if \eqref{eq:mumford-coniveau-input} holds at every
Hodge-generic point: countably many closed bounded-class loci cover
those points, so one dominates the complex curve.  Uniformly
bounded bridges on a dense set also suffice by first taking a
dominating Chow component and then arguing generically.
An isolated generic bridge is a pointwise input.  An isolated CM
bridge does not justify the generic endomorphism and Lie-action
identifications used in the theorem.
